% Results. Numbers: ../FINDINGS.md.
\section{Results}\label{sec:results}

\subsection{Exposure and outcomes}
Of the 10,495,850 candidates, 218,561 (2.1\%) cite at least one paper that had been retracted by $T_0$, and 172,149
are cited by one. Between February 2024 and October 2026, 4,323 candidates were retracted (0.041\%), 418 of them for
paper-mill reasons.

\subsection{Association and dose}
Candidates that cite a retracted paper were retracted later far more often (Table~\ref{tab:dose}). The rate rises from
0.035\% without a retracted reference to 0.27\% with one, 0.51\% with two and 0.83\% with three or four; the
1,451 candidates with five or more retracted references show 0.76\%, with a wide interval. The risk ratio for
$S_1 \ge 1$ is 8.8 (95\% CI 8.1--9.6). It is 17.8 (15.5--20.6) for citing a paper-mill retraction, and 6.8
(6.2--7.5) for being cited by a retracted paper.

\begin{table}[t]
\centering
\caption{Later retraction by number of references to papers retracted by $T_0$ (main analysis).}\label{tab:dose}
\begin{tabular}{lrrr}
\toprule
$S_1$ & Candidates & Retracted later & Rate (95\% CI) \\
\midrule
0 & 10,277,289 & 3,640 & 0.035\% (0.034--0.037) \\
1 & 187,735 & 502 & 0.27\% (0.25--0.29) \\
2 & 23,200 & 119 & 0.51\% (0.43--0.61) \\
3--4 & 6,175 & 51 & 0.83\% (0.63--1.08) \\
$\ge 5$ & 1,451 & 11 & 0.76\% (0.42--1.35) \\
\bottomrule
\end{tabular}
\end{table}

\subsection{Within journal and year}
Part of this association reflects journals: retractions cluster in journals and publication years, and papers in the
same journal cite one another. Within journal $\times$ publication year, the Mantel--Haenszel odds ratio for
$S_1 \ge 1$ is 2.99 (95\% CI 2.73--3.27; 1,201 informative strata holding 4,258 of the positives). By the
pre-registered rule, the signal holds within journal. It holds for paper-mill outcomes (3.37, 2.65--4.29) and for
other outcomes (2.92, 2.65--3.21). Being cited by a retracted paper gives a within-journal odds ratio of 3.76
(3.37--4.20). The estimate is unchanged without candidates whose DOI maps to several graph nodes (2.99) and slightly
higher for papers published 2015--2022 (3.10, 2.82--3.41).

\subsection{Screening value}
Despite the strong relative effect, the signal is weak as a screen (Table~\ref{tab:screen}). Most later retractions,
3,640 of 4,323 (84\%), cite no retracted paper at all. Among the 1,000 candidates with the most retracted references,
9 were retracted later (0.9\%, lift 22 over the base rate); among the top 10,000, 72 (0.72\%), which is 1.7\% of all
positives. The journal rate alone ranks one positive among its top 1,000 (0.1\%). The combined model, fitted on the earlier
window, has the highest average precision (0.0037) but a lower precision at 1,000 than the $S_1$ ranking.

\begin{table}[t]
\centering
\caption{Screening performance in the main analysis (95\% bootstrap intervals). Base rate 0.041\%.}\label{tab:screen}
\begin{tabular}{lrrrr}
\toprule
Ranking & Prec.@1,000 & Lift@1,000 & Prec.@10,000 & Avg. precision \\
\midrule
$S_1$ & 0.90\% (0.40--1.50) & 21.9 & 0.72\% (0.56--0.89) & 0.0013 \\
Journal rate $B$ & 0.10\% (0--0.30) & 2.4 & 0.02\% & 0.0023 \\
Combined model & 0.40\% (0.10--0.80) & 9.7 & 0.79\% (0.61--0.94) & 0.0037 \\
\bottomrule
\end{tabular}
\end{table}

\subsection{When the time point falls inside a retraction wave}
With $T_0' = $ 30 September 2023, the outcome window contains the mass retractions of late 2023 (7,578 positives,
3,141 of them paper-mill). The journal rate then becomes the best screen: 6.8\% of its top 1,000 candidates were
retracted (lift 94), against 1.1\% for $S_1$. The within-journal odds ratio for $S_1 \ge 1$ falls to 2.00
(1.86--2.15), and to 1.27 (1.13--1.42) for paper-mill outcomes, while it stays at 2.97 for other outcomes. By
January 2024 those waves had passed, and the journal rate placed only one future retraction among its top 1,000.

\subsection{Supporting analyses}
\paragraph{Community concentration.} The 27,567 retracted biomedical papers with a journal, their controls from the
same journal and year, and all their neighbours form a subgraph of 2.7 million papers and 57 million citations. Leiden
found 197 and 200 communities with two seeds (modularity 0.65 and 0.66); the two partitions agree only moderately
(adjusted Rand index 0.41), so both are reported. Retracted papers are more concentrated than their controls: the probability that
two retracted papers share a community is 1.15 and 1.11 times that of two controls (permutation $p = 0.001$ for
both, with a null mean of 1.00). The excess comes mainly from paper-mill retractions (1.68 and 1.52); for other
retractions it is small (1.06 and 1.05). Retracted paper-mill papers thus cluster in the citation graph beyond what
their journal and year imply, in line with \citet{feng2024molbio} and \citet{candal2024references}.

\paragraph{NIH exposure.} Of 2,332,727 NIH-funded biomedical papers published 2000--2023 in the graph, 10,872 (0.47\%)
cite at least one paper in a year after its retraction, compared with 0.48\% of other biomedical papers. Within journal
and year, NIH funding makes no practical difference (odds ratio 1.03, 1.00--1.05). NIH-funded papers cited 3,548
distinct retracted papers after their retraction.

\paragraph{Citations around retraction.} For 11,788 papers retracted 2000--2021 and matched controls from the same
journal and year, retracted papers received about 1.3 times the controls' citations in each of the five years before
the retraction year, 0.66 times in the year after and 0.47 times two years after. Citations continue after
retraction, at about two per paper in the second year, consistent with earlier studies
\citep{hsiao2021continued,schneider2020continued}.
